\documentclass[12pt]{article}
\usepackage[T1]{fontenc}
\usepackage[utf8]{inputenc}
\usepackage{lmodern}
\usepackage{textcomp}
\usepackage{enumitem}
\usepackage{geometry}
\usepackage{graphicx}
\usepackage{amsmath, amssymb}
\geometry{margin=1in}
\begin{document}
\begin{center}
Chemistry - Undergraduate Level Assessment 8 \
Total Marks: 40 \
Answer all questions.
\end{center}

\section*{Multiple Choice (10 marks)}
\begin{enumerate}[label=\arabic*.]
\item What is the maximum number of phases that can be at equilibrium with each other in a three component mixture?
\begin{enumerate}[label=(\alph*)]
    \item 2
    \item 3
    \item 4
    \item 5
\end{enumerate}
\item Which of the following is an n-type semiconductor?
\begin{enumerate}[label=(\alph*)]
    \item Silicon
    \item Diamond
    \item Silicon carbide
    \item Arsenic-doped silicon
\end{enumerate}
\item Considering 0.1 M aqueous solutions of each of the following, which solution has the lowest pH?
\begin{enumerate}[label=(\alph*)]
    \item Na$_{2}$CO$_{3}$
    \item Na$_{3}$PO$_{4}$
    \item Na$_{2}$S
    \item NaCl
\end{enumerate}
\item The fact that the infrared absorption frequency of deuterium chloride (DCl) is shifted from that of hydrogen chloride (HCl) is due to the differences in their
\begin{enumerate}[label=(\alph*)]
    \item electron distribution
    \item dipole moment
    \item force constant
    \item reduced mass
\end{enumerate}
\item What is the NMR frequency of 31 P in a 20.0 T magnetic field?
\begin{enumerate}[label=(\alph*)]
    \item 54.91 MHz
    \item 239.2 MHz
    \item 345.0 MHz
    \item 2167 MHz
\end{enumerate}
\end{enumerate}

\section*{True / False (10 marks)}
\textit{Answer True or False}
\begin{enumerate}[label=\arabic*.]
\setcounter{enumi}{5}
\item True or False: Each hybrid $sp^3$ orbital can hold a maximum of two electrons, not four.
\item True or False: The reduction of D-xylose with NaBH 4 produces a racemic mixture of enantiomers.
\item True or False: Calcium (Ca) has the lowest electron affinity among the elements listed due to its larger atomic size and lower effective nuclear charge.
\item True or False: The thermal stability of group-14 hydrides increases in the order: CH 4 \textless{} SiH 4 \textless{} GeH 4 \textless{} SnH 4 \textless{} PbH 4.
\item True or False: The polarization of a proton is 1.148 x 10^-6 at 0.335 T and 3.598 x 10^-5 at 10.5 T.
\end{enumerate}

\section*{Long Form Response (20 marks)}
\textit{Answer each question with a clear and complete explanation.}
\begin{enumerate}[label=\arabic*.]
\setcounter{enumi}{10}
\item Discuss how the Heisenberg uncertainty principle applies to a quantum mechanical particle in the lowest energy level of a one-dimensional box, focusing on the implications for position and momentum.
\item Discuss the decay of a radioactive isotope used in diagnostic imaging with a half-life of 6.0 hours. Calculate the remaining activity 24 hours post-delivery from an initial activity of 150 mCi, and explain the underlying principles of radioactive decay in your analysis.
\end{enumerate}

\vfill
\noindent\textit{End of Paper}
\end{document}